\section{Problema de negocio}

Artistas y organizaciones musicales deben decidir qué lanzamientos priorizar, cómo organizar un
catálogo y dónde concentrar recursos promocionales sin conocer completamente el desempeño futuro.
La pregunta central es:

\begin{quote}
¿Qué características musicales y contextuales están asociadas con la popularidad de una canción y
cómo puede esta información apoyar decisiones de promoción y gestión de catálogo?
\end{quote}

Los principales interesados y usos prudentes son:

\begin{itemize}
\item \textbf{Artistas:} comprender cómo se ubica una obra dentro de perfiles descriptivos, sin
convertirlos en reglas creativas.
\item \textbf{Sellos:} complementar revisión experta al priorizar experimentos o análisis de catálogo.
\item \textbf{Marketing musical:} caracterizar segmentos y formular campañas que luego deben medirse.
\item \textbf{Gestión de catálogo:} detectar grupos que requieren análisis específico y mantener una
visión documentada de cobertura.
\item \textbf{Curadores:} apoyar exploración de diversidad, nunca reemplazar juicio editorial.
\end{itemize}

Las restricciones son sustanciales: no hay fecha de extracción ni de lanzamiento, inversión,
exposición en playlists, región o comportamiento de usuarios. La popularidad es histórica y
dinámica según Spotify \citep{spotify_track_api}. En consecuencia, el análisis puede orientar
preguntas y priorizaciones, pero no garantizar desempeño comercial ni valorar mérito artístico.
